\section{Methodology}
\label{sec:methodology}

Building on the distinction in Section~\ref{sec:preliminary}, \method{} treats
semantic plans as the search objects and executable factors as the evaluation
objects. This separation is important for automated alpha discovery, where the
candidate space is large and market feedback is noisy and non-stationary.
Instead of searching directly over formulas or code, \method{} first explores
interpretable trading hypotheses and only realizes selected hypotheses into
executable factors.

The framework consists of four stages:
(1) semantic plan construction, which samples interpretable trading hypotheses
from a modular schema space; (2) plan-to-code realization, which converts
selected plans into executable factors with execution safeguards; (3)
reward-guided semantic search, which learns a surrogate model over evaluated
plans and allocates search budget among exploitation, exploration, and local
refinement; and (4) factor-pool consolidation, which filters validated factors
for downstream evaluation.


\subsection{Semantic Plan Space Construction}

\method{} searches over the structured plan space described in
Section~\ref{sec:preliminary}. The schema vocabulary can be constructed from
domain knowledge, existing factor libraries, and empirical discoveries from
previous runs. Unlike AST-based factor search, the schema specifies the economic
mechanism of a candidate factor while leaving its implementation flexible.

A semantic plan records the event, context, confirmation conditions,
directional hypothesis, and signal output format. This modular representation
creates a large combinatorial hypothesis space while preserving interpretability.
Moreover, the active schema vocabulary can be configured for different discovery
objectives without modifying the downstream realization and evaluation pipeline.


\subsection{Plan Realization and Execution Validation}

At each round, \method{} samples unseen plans from the active schema space and
selects a subset for implementation. For a selected plan \(p\), a code agent is
provided with the schema specification, data contract, and implementation
constraints, and generates a period-parameterized factor function
\(f_p(\cdot;\tau)\).

Before backtesting, generated factors are passed through execution guards that
verify contract compliance, computational validity, numerical stability, and
potential look-ahead leakage. Invalid implementations are removed before reward
computation, preventing erroneous evaluations from contaminating the search
process.


\subsection{Reward Modeling and Adaptive Semantic Search}

For search-time feedback, each valid realized factor receives a scalar reward
computed from its backtest statistics:
\begin{equation}
\begin{aligned}
r(p)=\frac{1}{4}\bigl(
&10\,\mathrm{IC}(p)+\mathrm{ICIR}(p) \\
&+10\,\mathrm{RankIC}(p)+\mathrm{RankICIR}(p)
\bigr),
\end{aligned}
\end{equation}
where \(\mathrm{IC}\) and \(\mathrm{RankIC}\) denote Pearson and rank information
coefficients, and \(\mathrm{ICIR}\) and \(\mathrm{RankICIR}\) are their
corresponding information ratios.

This reward is intentionally simple. It is used only to guide semantic search,
not as a learned utility function or as the final paper-level evaluation metric.
The multipliers on \(\mathrm{IC}\) and \(\mathrm{RankIC}\) are scale adjustments:
correlation coefficients are typically an order of magnitude smaller than their
information ratios, so rescaling prevents the averaged score from being dominated
by ratio terms. This design follows the same practical principle as
FactorEngine's composite fitness score over scaled predictive metrics
\cite{lin2026factorengine}: the search loop uses a transparent screening signal,
while final claims are reported with separate predictive and portfolio metrics.

All evaluated plans are stored in a reward buffer
\[
\mathcal{D}_t=\{(p_i,r_i,a_i)\}_{i=1}^{N_t},
\]
where \(a_i\) records execution status. Each plan is encoded into structured
schema features \(\phi(p)\). A LightGBM ensemble surrogate
\(g_{\theta_t}\) is trained on valid entries of \(\mathcal{D}_t\) to estimate the
expected reward of unseen plans:
\[
\hat r_t(p)=g_{\theta_t}(\phi(p)).
\]

For an unseen candidate \(p\), ensemble prediction provides a mean reward
estimate \(\mu_t(p)\) and uncertainty estimate \(\sigma_t(p)\). Candidate plans
are ranked using an upper confidence bound:
\[
\mathrm{UCB}_t(p)
=
\mu_t(p)+\kappa_t\sigma_t(p).
\]

Since selecting only the highest-UCB candidates may still cause premature
concentration, \method{} performs adaptive budget allocation:
\[
K=K_{\mu}+K_{\sigma}+K_{\mathrm{novel}},
\]
where the three components correspond to exploitation of high predicted reward
regions, uncertainty-driven exploration, and structural novelty over the schema
space. Diversity constraints are applied during batch construction to avoid
redundant evaluations.

In addition, \method{} maintains a local mutation channel that perturbs
historically high-reward plans through schema-level modifications, such as
changing contexts, adding or removing qualities, or replacing output forms.
This enables efficient refinement of promising semantic neighborhoods while the
main search process maintains global exploration.


\subsection{Validated Factor Pool}

After each search round, validated factors are added to the factor pool together
with their semantic plans, implementations, and evaluation statistics. A
diversity filter removes highly correlated or semantically redundant factors.
The final pool is obtained from the validated candidates and is used for
downstream portfolio evaluation.
